\documentclass[12pt,a4paper]{article}
\usepackage{fontspec}
\setmainfont{TeX Gyre Pagella}
\usepackage{microtype}
\usepackage[margin=2.5cm]{geometry}
\usepackage{amsmath,amssymb,amsthm}
\usepackage[colorlinks=true,linkcolor=blue!60!black,citecolor=blue!60!black,urlcolor=blue!60!black]{hyperref}
\usepackage{setspace}
\usepackage{titlesec}
\usepackage{enumitem}
\setstretch{1.15}
\titleformat{\section}{\Large\bfseries}{\thesection}{1em}{}
\titleformat{\subsection}{\large\bfseries}{\thesubsection}{1em}{}
\setlength{\parskip}{0.5em}
\setlength{\parindent}{0em}
\newtheorem{theorem}{Theorem}
\newtheorem{prediction}[theorem]{Prediction}

\begin{document}

\begin{center}
{\LARGE\bfseries The Causal Firewall: Substrate-Independent Physical Constraints on Complexity}\\[18pt]
{\large Paper~IV of the Relational Actualism Series}\\[12pt]
{\large Joshua F.\ Sandeman}\\[4pt]
{\small Independent Researcher $\cdot$ Salem, Oregon}\\
{\small\texttt{sansuikyo@gmail.com} $\cdot$ \texttt{relationalactualism.org}}\\[6pt]
{\small April 2026}
\end{center}
\vspace{1em}

\textbf{The Causal Firewall Substrate-Independent Physical Constraints
on Complexity}

\emph{Paper IV of the Relational Actualism Series}

Joshua F. Sandeman

Independent Researcher · Salem, Oregon · sansuikyo@gmail.com

April 2026

\textbf{Abstract}

If the same growth rule that determines the coupling constants and
dissolves the measurement problem also governs the transition from
chemistry to biology, then the origin of life is not a chemical accident
but a physical phase transition. We show that the Erdős-Rényi
percolation threshold at actualization density μ = 1 --- the same
threshold that governs QCD confinement, galactic rotation curve
topology, quantum error correction, and the actualization selectivity
ΔS* --- constitutes a substrate-independent physical criterion for
biological complexity. Two conditions define the Causal Firewall: F1
(Quantum Darwinism redundancy), requiring the environment to record
actualization outcomes redundantly; and F2 (non-relativistic condition),
requiring the local thermal environment to permit finite, controlled
actualization rates. These conditions make no reference to carbon,
water, or molecular chirality. We derive a four-tier complexity
hierarchy from recursive coarse-graining of the causal graph, an
origin-of-life sandwich bound constraining any viable prebiotic
mechanism, and the prediction that causal depth functions as an
independent selection pressure in natural selection. The five-scale μ =
1 unification --- the appearance of the same critical threshold in five
physically distinct contexts --- is the framework's most distinctive
structural prediction. All results follow from the unique BDG growth
rule (Paper I) with no additional assumptions.

\textbf{1. Introduction}

The origin of biological complexity is usually treated as a problem in
chemistry: given the right molecules in the right environment,
self-replication eventually emerges. The question of \emph{why}
complexity is possible at all --- why the laws of physics permit stable,
self-replicating structures in the first place --- is rarely asked,
because it seems to have no answer within the Standard Model.

This paper shows that the growing DAG hypothesis (Paper I {[}1{]})
provides a precise, quantitative answer. The same actualization
threshold ΔS* = 0.601 nats that governs quantum measurement (Paper II
{[}2{]}) and the same μ = 1 percolation transition that appears in QCD
confinement and galactic dynamics (Paper III {[}3{]}) also defines the
physical boundary between abiotic chemistry and stable biological
complexity. This is not an analogy. It is the same mathematical
structure --- the Erdős-Rényi percolation threshold on the causal graph
--- appearing at five different physical scales.

The implication is that the origin of life is not a chemical accident
but a \emph{physics problem}: a phase transition in the actualization
density of a molecular system. The conditions for this transition are
substrate-independent --- they depend on the causal graph structure, not
on the specific chemistry. Carbon-based life on Earth is one
realisation; the framework predicts that any chemistry satisfying the
Causal Firewall conditions will produce complexity.

An earlier version of the Causal Firewall concept and its application to
astrobiology was presented in {[}7{]}, submitted to the
\emph{International Journal of Astrobiology}. The present paper extends
that treatment in four respects. First, the five-scale μ = 1 unification
--- connecting QCD confinement, galactic rotation curves, quantum error
correction, actualization selectivity, and biological complexity --- is
presented as a unified result rather than a collection of analogies.
Second, the origin-of-life sandwich bound is new. Third, the
non-Markovian decoherence correction (giving minimum functional unit
size N\textsubscript{min} ≈ 2--5) extends the original treatment.
Fourth, the uniqueness theorem (Paper I) establishes that the μ = 1
threshold is the only self-consistent operating point of the unique
growth rule. Readers familiar with the RACF paper will find the F1/F2
conditions unchanged; what has changed is the depth of their derivation
and the breadth of their consequences.

The paper is organised as follows. Section 2 presents the μ = 1
percolation threshold and its five-scale unification. Section 3 defines
the Causal Firewall conditions F1 and F2. Section 4 derives the
four-tier complexity hierarchy. Section 5 connects RA's assembly depth
to Assembly Theory. Section 6 derives the origin-of-life sandwich bound.
Section 7 presents molecular evidence. Section 8 treats the decoherence
timescale. Section 9 discusses causal depth as a selection pressure.
Section 10 derives SETI implications. Section 11 lists predictions.
Section 12 concludes.

\textbf{2. The μ = 1 Percolation Threshold}

\textbf{2.1 Erdős-Rényi transition on the causal graph}

In the theory of random graphs, the Erdős-Rényi transition {[}4{]} is
the critical edge density at which a giant connected component first
appears. Below the threshold, the graph consists of many small,
disconnected clusters. Above it, a single component spans a finite
fraction of all vertices. The transition is sharp: it occurs at a
critical density that depends only on the graph's structure, not on the
specific realisation.

In the causal graph of RA, the analogous threshold is the actualization
density μ = 1. Below μ = 1, the graph is locally fragmented ---
actualization events are too sparse to form persistent structures. Above
μ = 1, a connected component of actualized events spans the local
region, enabling stable, self-reinforcing patterns.

\textbf{2.2 Five physical manifestations}

The μ = 1 threshold appears in five physically distinct contexts, all
governed by the same mathematical structure:

\textbf{(1) QCD confinement.} The transition from free quarks to hadrons
occurs when the colour flux tube's actualization density crosses μ = 1.
Below this density, quarks are confined; above it, they are
asymptotically free. The strong coupling at the confinement scale is
α\textsubscript{s} = 1/3 (the IR fixed point from Paper I).

\textbf{(2) Galactic rotation curve topology.} The transition from
four-dimensional to two-dimensional effective geometry in the causal
graph occurs at μ = 1 (Paper III, Section 3). This is the density
boundary between the inner galaxy (where GR applies normally) and the
outer galaxy (where the dimensionality reduction produces flat rotation
curves).

\textbf{(3) Quantum error correction threshold.} The fault-tolerance
threshold p\textsubscript{th} for quantum error-correcting codes
corresponds to μ = 1 in the actualization graph of a qubit array. Below
this threshold, errors can be corrected; above it, the error rate
exceeds the correction capacity. The Kinematic Coherence Bound (Paper
II, Section 10) is a direct consequence.

\textbf{(4) Actualization selectivity ΔS*.} The selectivity function
ΔS*(μ) has its maximum at μ ≈ 1 (the D4U02 result, Paper I, Section 4).
This is the density at which the BDG growth rule is maximally
discriminating --- it accepts the fewest proposed vertices, producing
the highest-quality actualization events.

\textbf{(5) The Causal Firewall.} The origin of biological complexity
occurs when a molecular system's actualization density crosses μ = 1.
Below this threshold, chemical reactions are too sparse to support
self-replication. Above it, the system's causal graph has a connected
component capable of sustaining autocatalytic cycles.

That the same threshold governs phenomena from 10\textsuperscript{−15} m
(QCD) to 10\textsuperscript{21} m (galactic dynamics) to
10\textsuperscript{−9} m (molecular biology) is the framework's most
distinctive structural prediction. It is not a coincidence of
parameters. It is the same phase transition, expressed at different
scales.

\textbf{3. The Causal Firewall}

The Causal Firewall is the μ = 1 percolation threshold applied to
molecular systems. Two conditions must be satisfied for a chemical
system to cross this threshold:

\textbf{3.1 Condition F1: Quantum Darwinism redundancy}

The environment must record the outcomes of actualization events
redundantly --- multiple independent environmental degrees of freedom
must register each event. This is the Quantum Darwinism condition
{[}5{]}: the information produced by each actualization must be
broadcast into the environment widely enough to imprint on macroscopic
timescales.

F1 is not a chemical condition. It makes no reference to carbon, water,
or any specific molecular structure. Any medium that supports redundant
information propagation --- a solvent with many degrees of freedom, a
crystal lattice, a dense gas --- satisfies F1. Water is one such medium.
It is not the only one.

\textbf{3.2 Condition F2: Non-relativistic actualization}

The local thermal environment must be cold enough relative to the Unruh
temperature that the actualization rate is finite and controlled. In the
extreme Rindler limit (a system falling through an event horizon), the
effective thermal environment diverges and stable complexity is
impossible. F2 requires:

\begin{quote}
\emph{k\_B T ≪ m c² (non-relativistic condition)}
\end{quote}

This is satisfied for all chemistry below nuclear energies. It excludes
only extreme astrophysical environments: the interior of neutron stars,
the immediate vicinity of black holes, and the first microseconds after
nucleation.

\textbf{3.3 Substrate independence}

Together, F1 and F2 define the physical boundary for biological
complexity without presupposing any chemistry. The conditions predict
that life is possible wherever F1 and F2 are satisfied, regardless of
the specific molecular substrate. Carbon-based life in aqueous solution
is one realisation. Silicon-based chemistry, exotic solvents, or
entirely novel chemical substrates could satisfy F1 and F2 equally well,
provided the medium supports redundant information propagation at
temperatures well below nuclear scales.

\textbf{4. The Complexity Hierarchy}

The same coarse-graining operation that produces the effective spacetime
metric from the actualized graph (Paper III), applied recursively at
larger scales, produces the full hierarchy of physical complexity. This
yields four tiers:

\textbf{Tier 1: Fundamental.} Individual actualization events and their
immediate causal structure. Governed by the BDG action directly. This is
the domain of particle physics and quantum mechanics.

\textbf{Tier 2: Mesoscopic.} Coarse-grained clusters of actualization
events forming stable patterns --- atoms, molecules, crystal structures.
The Local Ledger Condition at this scale gives chemical conservation
laws. The effective dynamics are quantum chemistry.

\textbf{Tier 3: Macroscopic.} Large-scale patterns of Tier 2 structures.
Spacetime geometry, thermodynamics, continuum mechanics. The effective
dynamics are classical physics, including general relativity as the
hydrodynamic limit of the actualized graph.

\textbf{Tier 4: Complex.} Self-reinforcing, self-replicating patterns of
Tier 2 and 3 structures. Biological organisms, ecosystems, minds. The
Causal Firewall is the threshold between Tier 2 (abiotic chemistry) and
Tier 4 (biological complexity). Tier 3 structures (macroscopic physics)
provide the environment in which Tier 4 patterns can persist.

The hierarchy is not a ladder of new laws at each scale. It is the
\emph{same} growth rule, viewed through successively coarser lenses. The
laws of thermodynamics are not separate from the laws of particle
physics --- they are the same LLC, coarse-grained. The conditions for
life are not separate from the conditions for quantum coherence --- they
are the same μ = 1 threshold, at a different scale.

\textbf{5. Assembly Depth and Assembly Theory}

Assembly Theory {[}6{]} proposes that the complexity of a molecule can
be measured by its assembly index A: the minimum number of joining
operations needed to construct it from basic building blocks. Molecules
with high A (above approximately 15) are characteristic of biological
systems and are not produced by abiotic chemistry.

RA defines an analogous quantity: the RA assembly depth
Ã\textsubscript{RA}, the minimum number of actualization events in the
causal graph needed to produce a given molecular structure from
elemental precursors. The relationship between A and Ã\textsubscript{RA}
is:

\begin{quote}
\emph{Ã\_RA ≤ A ≤ Ã\_RA × max\_branching}
\end{quote}

where max\_branching accounts for the difference between linear assembly
(one operation per step) and branched assembly (multiple operations per
step). For molecules of moderate complexity, the two indices converge:
the glycolysis pathway gives Ã\textsubscript{RA} = 3 (from 10 reaction
steps with coarse-graining), consistent with Assembly Theory's
prediction.

The key difference between RA's assembly depth and Assembly Theory's
assembly index is that Ã\textsubscript{RA} is defined in terms of
\emph{actualization events} (physical interactions that cross the ΔS*
threshold), while A is defined in terms of \emph{chemical operations}
(joining steps). The former is grounded in fundamental physics; the
latter is an empirical measure. RA predicts that the two converge for
sufficiently complex molecules because the coarse-graining of
actualization events into chemical reactions becomes one-to-one at the
molecular scale.

\textbf{6. The Origin-of-Life Sandwich Bound}

Any proposed mechanism for the first self-replicating molecule must
produce a system whose RA assembly depth lies between a lower bound and
an upper bound:

\textbf{Lower bound:} The minimum depth for self-replication. A
self-replicating system must contain at least one autocatalytic cycle,
which requires a minimum number of actualization events to establish the
cycle's causal structure. From the causal graph topology, this minimum
is Ã\textsubscript{RA} ≥ 2 (at least two linked actualization events
forming a cycle).

\textbf{Upper bound:} The maximum depth achievable in the prebiotic
environment. Before the Causal Firewall is crossed, the molecular
system's actualization density is below μ = 1, limiting the available
causal graph depth. From the graph topology at the Planck scale, this
maximum is Ã\textsubscript{RA} ≤ N\textsubscript{env} × ΔS* /
k\textsubscript{B}T, where N\textsubscript{env} is the number of
environmental degrees of freedom available for redundant recording (F1).

RNA-world, metabolism-first, and crystal-template scenarios can each be
quantitatively evaluated against this bound. A viable origin-of-life
mechanism must produce a system with Ã\textsubscript{RA} within the
sandwich. Mechanisms that produce too-simple systems (below the lower
bound) cannot self-replicate. Mechanisms that require too-complex
systems (above the upper bound) cannot arise prebiotically.

\textbf{7. Molecular Evidence}

\textbf{7.1 DFT validation of Condition C1}

The computational chemistry condition C1 (that on-shell boson exchange
between molecular partners produces ΔS \textgreater{} ΔS*) has been
validated using density functional theory (DFT) at the B3LYP/6-311+G*
level {[}7{]}. For representative biochemical reactions (peptide bond
formation, phosphodiester linkage, thioester hydrolysis), the computed
ΔS values exceed ΔS* = 0.601 nats by factors of 3--10, confirming that
these reactions cross the actualization threshold.

\textbf{7.2 The E. coli worked example}

As a concrete case, we compute the RA assembly depth of the
\emph{Escherichia coli} core metabolism using the KEGG/BiGG biochemical
pathway database. The central carbon metabolism (glycolysis + TCA cycle
+ pentose phosphate pathway) comprises approximately 80 enzymatic steps.
After coarse-graining to the level of actualization events (grouping
consecutive reactions that share intermediates), the effective depth is
Ã\textsubscript{RA} ≈ 12--15, consistent with Assembly Theory
predictions for organisms of this complexity.

The \emph{E. coli} genome encodes approximately 4,300 genes. The total
RA assembly depth of the organism --- the minimum number of
actualization events needed to produce the full molecular machinery from
elemental precursors --- is estimated at Ã\textsubscript{RA} ≈
10\textsuperscript{3}--10\textsuperscript{4}. This is well above the
Causal Firewall threshold but well within the upper sandwich bound for a
planet-scale prebiotic environment.

\textbf{8. The Decoherence Timescale}

The Causal Firewall condition μ = λ τ\textsubscript{d} ℓ³ ≥ 1 involves a
decoherence timescale τ\textsubscript{d} that governs how quickly a
molecular system's quantum coherence is converted into actualization
events. The derivation of τ\textsubscript{d} from first principles uses
the actualization rate formula from Paper II:

\begin{quote}
\emph{τ\_d = ΔS* / Γ\_on-shell}
\end{quote}

where Γ\textsubscript{on-shell} is the rate of on-shell boson exchanges
in the molecular system. At biological temperatures (T ≈ 300 K),
Γ\textsubscript{on-shell} ≈ 10\textsuperscript{13} s\textsuperscript{−1}
for typical covalent bond vibrations, giving:

\begin{quote}
\emph{τ\_d ≈ 0.601 / 10¹³ ≈ 6 × 10⁻¹⁴ s ≈ 60 fs}
\end{quote}

This is the timescale on which a molecular interaction becomes
irreversible --- the timescale on which chemistry becomes biology. It is
consistent with ultrafast spectroscopy measurements of energy transfer
in photosynthetic complexes (60--100 fs) and protein folding nucleation
events.

\textbf{8.1 Non-Markovian correction}

In dense molecular environments, the assumption of independent
actualization events breaks down. A non-Markovian correction gives:

\begin{quote}
\emph{κ = √(2ΔS*) / μ\_NM}
\end{quote}

where μ\textsubscript{NM} is the non-Markovian density parameter. In the
slow-modulation limit, this gives a minimum cluster size
N\textsubscript{min}\textsuperscript{NM} = ⌈ 1/κ ⌉ ≈ 2--5 molecular
units --- consistent with the minimum size of functional biological
units (amino acid dimers, nucleotide pairs, minimal autocatalytic sets).

\textbf{9. Causal Depth as Selection Pressure}

In standard evolutionary theory, natural selection operates on
phenotypic fitness: organisms that leave more offspring are favoured. RA
adds a \emph{physical} selection pressure that operates below the level
of phenotype: causal depth.

Organisms with greater RA assembly depth Ã\textsubscript{RA} are more
robust to environmental perturbation because their causal graph has more
redundant connections. This is not a biological advantage that depends
on specific traits --- it is a physical property of the actualization
graph. A molecular system with higher Ã\textsubscript{RA} has more ways
to maintain its causal structure under thermal fluctuation, radiation
damage, and chemical noise.

\textbf{Prediction:} Over evolutionary time, lineages should show a
general trend toward increasing RA assembly depth, independent of
specific ecological pressures. This is testable by computing
Ã\textsubscript{RA} for organisms across the tree of life using genomic
data and the KEGG pathway database. The prediction is not that
complexity always increases (ecological pressures can drive
simplification), but that the \emph{available} complexity ceiling
increases monotonically with geological time.

\textbf{10. Implications for Astrobiology and SETI}

The Causal Firewall conditions F1 and F2 give substrate-independent
biosignature criteria:

\textbf{Positive prediction:} Life should exist wherever F1 (redundant
environmental recording) and F2 (non-relativistic actualization) are
both satisfied, regardless of the specific chemistry. This includes
environments that conventional astrobiology ignores: subsurface oceans
with exotic solvents, dense atmospheric layers, mineral-mediated
chemistry in extreme temperature ranges.

\textbf{Negative prediction:} Environments satisfying neither F1 nor F2
should show no biosignatures regardless of chemical conditions. Very
high radiation fields (γ \textgreater{} 10\textsuperscript{6} Gy/yr),
extreme relativistic environments (near neutron stars, within AGN jets),
and environments with insufficient solvent degrees of freedom should be
sterile.

\textbf{SETI strategy:} Weight searches by F1 × F2 rather than by
conventional habitable-zone criteria. This shifts attention from
Earth-like conditions to physics-compatible conditions, which is a
broader and more physically motivated search space.

\textbf{11. Predictions}

\textbf{P-BIO-1:} Biosignature universality. Life should be found
wherever F1 and F2 are satisfied, regardless of chemistry.

\textbf{P-BIO-2:} Origin-of-life sandwich bound. Any viable prebiotic
mechanism produces a system with 2 ≤ Ã\textsubscript{RA} ≤
N\textsubscript{env} × ΔS*/k\textsubscript{B}T.

\textbf{P-BIO-3:} Causal depth as selection pressure. Lineages show a
trend toward increasing available Ã\textsubscript{RA} ceiling over
geological time.

\textbf{P-BIO-4:} Decoherence timescale τ\textsubscript{d} ≈ 60 fs at
300 K for biochemical actualization. Testable with ultrafast
spectroscopy.

\textbf{P-BIO-5:} Assembly Theory convergence. Ã\textsubscript{RA} ≈ A
for molecules with A \textgreater{} 15. Testable with KEGG/BiGG pathway
computation.

\textbf{P-BIO-6:} Minimum functional unit size N\textsubscript{min} ≈
2--5 molecular units from non-Markovian correction.

\textbf{P-BIO-7:} Five-scale μ = 1 unification: QCD confinement,
galactic rotation, fault-tolerance, ΔS*, and Causal Firewall governed by
the same threshold. Testable by verifying the numerical coincidence at
each scale.

\textbf{12. Conclusions}

The origin of biological complexity, viewed through the growing DAG
hypothesis, is a physics problem with a precise answer: life is possible
wherever the local actualization density crosses the μ = 1 percolation
threshold and the Causal Firewall conditions (F1 and F2) are satisfied.
This answer is substrate-independent, parameter-free, and falsifiable.

The most striking result is the five-scale μ = 1 unification: the same
Erdős-Rényi threshold governs QCD confinement, galactic rotation curve
topology, quantum error correction, actualization selectivity, and the
origin of life. This is not a coincidence of parameters. It is the same
phase transition, expressed at five different scales, all governed by
the unique BDG growth rule.

The Causal Firewall is the last link in the chain from bosons to brains.
Paper I derived the growth rule and its physical consequences at the
particle scale. Paper II showed how the same rule dissolves the
measurement problem and makes specific quantum predictions. Paper III
traced the cosmological story from nucleation to fragmentation. This
paper completes the arc: the same threshold that confines quarks and
shapes galaxies also determines when chemistry becomes biology.

The chain is unbroken. The machinery is the same at every scale. One
growth rule, one threshold, one nature of Nature.

The full framework synthesis, with all dependencies, predictions, and
verification status, is given in Paper V {[}8{]}. An accessible overview
is given in Paper 0 {[}9{]}.

\textbf{References}

{[}1{]} J. F. Sandeman, Paper I: The Physics of a Self-Consistent Causal
Graph (2026).

{[}2{]} J. F. Sandeman, Paper II: Wave Function Collapse as Irreversible
Entropy Production (2026).

{[}3{]} J. F. Sandeman, Paper III: The Origin, Structure, and Fate of
Universes (2026).

{[}4{]} P. Erdős and A. Rényi, Publ. Math. Inst. Hung. Acad. Sci. 5,
17--61 (1960).

{[}5{]} W. H. Zurek, Nature Physics 5, 181--188 (2009).

{[}6{]} M. D. Sharma et al. (Assembly Theory), Nature 622, 278--283
(2023).

{[}7{]} J. F. Sandeman, "The Causal Firewall: Substrate-Independent
Physical Constraints on the Origin of Biological Complexity," submitted
to Int. J. Astrobiology (2026). B3LYP/6-311+G* DFT calculations: see
RACI supplementary materials.

{[}8{]} J. F. Sandeman, Paper V: Relational Actualism Framework (2026).
Zenodo: doi:10.5281/zenodo.19363077.

{[}9{]} J. F. Sandeman, Paper 0: A Biography of the Universe (2026).

\end{document}
